\begin{table}[H]
\centering
\caption{Standardized Effect Sizes (SDE) Appendix}
\label{tab:sde}
\footnotesize
\begin{tabular}{lccccc}
\hline\hline
& $\hat{\beta}$ & SE & SD($Y$) & SDE & Classification \\
\hline
\multicolumn{6}{c}{\textit{Panel A: Pooled Effect}} \\
\hline
All sectors (10\% MW increase) & $-0.0042$ & $0.0016$ & 0.0187 & $-0.224$ & Small Negative \\
\hline
\multicolumn{6}{c}{\textit{Panel B: Heterogeneous Effects}} \\
\hline
High-Kaitz (0.50+), 10\% MW & $-0.0042$ & $0.0016$ & 0.0187 & $-0.224$ & Small Negative \\
Low-Kaitz (0.20-), 10\% MW & $-0.0017$ & $0.0014$ & 0.0187 & $-0.091$ & Small Negative \\
Ratio (High:Low) & & & & 2.46$\times$ & \\
\hline\hline
\end{tabular}

\begin{flushleft}
\small

\textit{Notes:}

\textbf{Country:} Slovakia (EU member state, Central Europe).

\textbf{Research question:} Does a sustained escalation of the statutory minimum wage (129\% nominal increase 2012--2024) reduce employment differentially across sectors depending on wage-floor bite (Kaitz index)?

\textbf{Policy mechanism:} Government decree under Act 663/2007 sets a uniform national gross monthly minimum wage that applies to all sectors simultaneously. No sector carve-outs or exemptions. When the wage floor rises sharply, sectors where the floor represents a larger fraction of median wages (high Kaitz) experience a larger shock to marginal labor costs than low-Kaitz sectors.

\textbf{Outcome definition:} Log quarterly employment by sector from Eurostat Labour Force Survey (\texttt{lfsq\_egan2}), aggregated into 22 NACE 2-digit sectors. Outcome measured in thousands of employed persons per quarter.

\textbf{Treatment:} Pre-determined 2012 Kaitz index (statutory minimum wage / sector median wage) interacted with staggered log minimum wage increases 2013--2024. High-Kaitz sectors (Accommodation, Retail: 0.38--0.61) vs. low-Kaitz sectors (Finance, ICT: 0.19--0.20).

\textbf{Data:} Quarterly employment panel from Eurostat 2010-Q1 through 2025-Q4 (64 quarters), 22 sectors, 1,408 sector-quarter observations. Wage data from Eurostat Labour Cost Index. Statutory minimum wages from Slovak Ministry of Labor. All sources are public EU statistics with no authentication required.

\textbf{Method:} Difference-in-differences with continuous treatment (Kaitz index $\times$ cumulative log MW change), estimated by OLS with sector and quarter fixed effects. Standard errors clustered by sector. Specification controls for pre-2012 sector employment level and prior employment growth. Event study validates no anticipation or pre-trends.

\textbf{Sample:} All 22 NACE sectors with continuous employment data 2010--2025. No sector exclusions or sample restrictions. Balanced panel across all time periods. High-Kaitz subset includes 10 sectors with index $\geq 0.35$; low-Kaitz subset includes 12 sectors with index $< 0.25$.

SDE $= \hat{\beta} / \text{SD}(Y)$, where $\hat{\beta}$ is the estimated coefficient (log-point change in employment) and SD$(Y)$ is the standard deviation of quarterly log employment changes. For high-Kaitz sectors experiencing the mean 10\% minimum wage increase, SDE is calculated as: $(-0.084 \times 0.50 \times 0.10) / 0.0187 = -0.224$. For low-Kaitz sectors: $(-0.084 \times 0.20 \times 0.10) / 0.0187 = -0.091$.

Classification refers to magnitude, not statistical significance: Large ($|SDE| > 0.15$), Moderate ($0.05 < |SDE| \leq 0.15$), Small ($0.005 < |SDE| \leq 0.05$), Null ($|SDE| \leq 0.005$).

\end{flushleft}
\end{table}
